\documentclass[10pt,a4paper]{amsart}
\usepackage{amsmath,amssymb,amsthm,graphicx}
% Retained minimal preamble for the Stacks Project chapter more-morphisms.tex (More on Morphisms).
% Verbatim excerpt of the chapter preamble: Xy-pic setup (native preamble lines 14-18) and the native
% theorem declarations with the mathematical macros (native preamble lines 154-217). Copied mechanically;
% nothing added, removed or redefined. The wrapper class and the amsmath/amssymb/amsthm/graphicx base
% packages are supplied by the assembler, so the class/externaldocument part of the native preamble is omitted.
% For commutative diagrams we use Xy-pic
\usepackage[all]{xy}

% We use 2cell for 2-commutative diagrams.
\xyoption{2cell}
% Theorem environments.
%
\theoremstyle{plain}
\newtheorem{theorem}[subsection]{Theorem}
\newtheorem{proposition}[subsection]{Proposition}
\newtheorem{lemma}[subsection]{Lemma}

\theoremstyle{definition}
\newtheorem{definition}[subsection]{Definition}
\newtheorem{example}[subsection]{Example}
\newtheorem{exercise}[subsection]{Exercise}
\newtheorem{situation}[subsection]{Situation}

\theoremstyle{remark}
\newtheorem{remark}[subsection]{Remark}
\newtheorem{remarks}[subsection]{Remarks}

\numberwithin{equation}{subsection}

% Macros
%
\def\lim{\mathop{\mathrm{lim}}\nolimits}
\def\colim{\mathop{\mathrm{colim}}\nolimits}
\def\Spec{\mathop{\mathrm{Spec}}}
\def\Hom{\mathop{\mathrm{Hom}}\nolimits}
\def\Ext{\mathop{\mathrm{Ext}}\nolimits}
\def\SheafHom{\mathop{\mathcal{H}\!\mathit{om}}\nolimits}
\def\SheafExt{\mathop{\mathcal{E}\!\mathit{xt}}\nolimits}
\def\Sch{\mathit{Sch}}
\def\Mor{\mathop{\mathrm{Mor}}\nolimits}
\def\Ob{\mathop{\mathrm{Ob}}\nolimits}
\def\Sh{\mathop{\mathit{Sh}}\nolimits}
\def\NL{\mathop{N\!L}\nolimits}
\def\CH{\mathop{\mathrm{CH}}\nolimits}
\def\proetale{{pro\text{-}\acute{e}tale}}
\def\etale{{\acute{e}tale}}
\def\QCoh{\mathit{QCoh}}
\def\Ker{\mathop{\mathrm{Ker}}}
\def\Im{\mathop{\mathrm{Im}}}
\def\Coker{\mathop{\mathrm{Coker}}}
\def\Coim{\mathop{\mathrm{Coim}}}

% Boxtimes
%
\DeclareMathSymbol{\boxtimes}{\mathbin}{AMSa}{"02}

%
% Macros for moduli stacks/spaces
%
\def\QCohstack{\mathcal{QC}\!\mathit{oh}}
\def\Cohstack{\mathcal{C}\!\mathit{oh}}
\def\Spacesstack{\mathcal{S}\!\mathit{paces}}
\def\Quotfunctor{\mathrm{Quot}}
\def\Hilbfunctor{\mathrm{Hilb}}
\def\Curvesstack{\mathcal{C}\!\mathit{urves}}
\def\Polarizedstack{\mathcal{P}\!\mathit{olarized}}
\def\Complexesstack{\mathcal{C}\!\mathit{omplexes}}
% \Pic is the operator that assigns to X its picard group, usage \Pic(X)
% \Picardstack_{X/B} denotes the Picard stack of X over B
% \Picardfunctor_{X/B} denotes the Picard functor of X over B
\def\Pic{\mathop{\mathrm{Pic}}\nolimits}
\def\Picardstack{\mathcal{P}\!\mathit{ic}}
\def\Picardfunctor{\mathrm{Pic}}
\def\Deformationcategory{\mathcal{D}\!\mathit{ef}}
\newenvironment{reference}{}{}
\newenvironment{slogan}{}{}
\newenvironment{history}{}{}
\usepackage{cleveref}
\crefname{originallemma}{original lemma}{original lemmas}
\Crefname{originallemma}{Original lemma}{Original lemmas}
\makeatletter
\newlabel{more-algebra-proposition-finite-type-over-G-ring}{{07PV}{source}{}{sourcepointer-1}{}}
\newlabel{more-algebra-proposition-finite-type-over-G-ring@cref}{{[originallemma][0][]07PV}{[1][0][]source}{}{}{}}
\makeatother
\crefname{originallemma}{original lemma}{original lemmas}
\Crefname{originallemma}{Original lemma}{Original lemmas}
\makeatletter
\newlabel{smoothing-lemma-approximation-property-variant}{{0CAR}{source}{}{sourcepointer-2}{}}
\newlabel{smoothing-lemma-approximation-property-variant@cref}{{[originallemma][0][]0CAR}{[1][0][]source}{}{}{}}
\makeatother
\title{Artin approximation of a local map from a scheme into another by an etale neighbourhood}
\author{The Stacks Project Authors}
\date{stacks-project:more-morphisms.tex:local-snapshot-20261011; GNU FDL 1.2 or later}
\begin{document}
\maketitle
\noindent\textit{Editorial note.} Statement span is lines 11016-11039 of sources/stacks-more-morphisms/pinned-original-source.tex; proof span is lines 11041-11068, the complete author proof (26 body lines).\par\medskip
\noindent\textit{Editorial note.} Navigation stubs: 2 cited-but-not-included labels are supplied as non-load-bearing pointers with their pinned tags where the snapshot tag index has them.\par\medskip
\noindent\textit{Editorial note.} extra\_packages is empty on purpose: the retained preamble already loads xy with options and a second declaration clashes in the assembler wrapper.\par\medskip
\noindent\textit{Editorial note.} Text is retained verbatim, including the author's own typos and placeholders.\par\medskip
\noindent\textit{Allowed standard prerequisites (editorial).}\begin{itemize}
\item Explicit prerequisite boundary. G-rings, completions of local rings, elementary etale neighbourhoods, and Artin approximation; the construction is retained.
\end{itemize}\medskip
\ifdefined\hypertarget\hypertarget{sourcepointer-1}{}\fi
\noindent\textit{Source reference.} Original lemma 07PV: sources/stacks-more-morphisms/pinned-original-source.tex, source lines 11068--11068. Cited by the retained statement or proof but not part of the collected spans; supplied only as a navigation pointer. The argument does not depend on its content.\par\medskip
\ifdefined\hypertarget\hypertarget{sourcepointer-2}{}\fi
\noindent\textit{Source reference.} Original lemma 0CAR: sources/stacks-more-morphisms/pinned-original-source.tex, source lines 11068--11068. Cited by the retained statement or proof but not part of the collected spans; supplied only as a navigation pointer. The argument does not depend on its content.\par\medskip
\begin{lemma}
\label{lemma-map-approximation}
Let $S$ be a locally Noetherian scheme. Let $X$, $Y$ be
schemes locally of finite
type over $S$. Let $x \in X$ and $y \in Y$ be points lying over the
same point $s \in S$. Assume $\mathcal{O}_{S, s}$ is a G-ring.
Assume further we are given a local $\mathcal{O}_{S, s}$-algebra map
$$
\varphi : \mathcal{O}_{Y, y} \longrightarrow \mathcal{O}_{X, x}^\wedge
$$
For every $N \geq 1$
there exists an elementary \'etale neighbourhood
$(U, u) \to (X, x)$ and an $S$-morphism
$f : U \to Y$ mapping $u$ to $y$ such that the diagram
$$
\xymatrix{
\mathcal{O}_{X, x}^\wedge \ar[r] &
\mathcal{O}_{U, u}^\wedge \\
\mathcal{O}_{Y, y} \ar[r]^{f^\sharp_u} \ar[u]^\varphi &
\mathcal{O}_{U, u} \ar[u]
}
$$
commutes modulo $\mathfrak m_u^N$.
\end{lemma}

\begin{proof}
The question is local on $X$ hence we may assume $X$, $Y$, $S$ are affine.
Say $S = \Spec(R)$, $X = \Spec(A)$, $Y = \Spec(B)$.
Write $B = R[x_1, \ldots, x_n]/(f_1, \ldots, f_m)$.
Let $\mathfrak p \subset A$ be the prime ideal corresponding to $x$.
The local ring $\mathcal{O}_{X, x} = A_\mathfrak p$ is a G-ring by
More on Algebra, Proposition
\ref{more-algebra-proposition-finite-type-over-G-ring}.
The map $\varphi$ is a map
$$
B_\mathfrak q^\wedge \longrightarrow A_\mathfrak p^\wedge
$$
where $\mathfrak q \subset B$ is the prime corresponding to $y$.
Let $a_1, \ldots, a_n \in A_\mathfrak p^\wedge$ be the images
of $x_1, \ldots, x_n$ via
$R[x_1, \ldots, x_n] \to B \to B_\mathfrak q^\wedge \to A_\mathfrak p^\wedge$.
Then we can apply Smoothing Ring Maps, Lemma
\ref{smoothing-lemma-approximation-property-variant}
to get an \'etale ring map $A \to A'$ and a prime ideal
$\mathfrak p' \subset A'$ and $b_1, \ldots, b_n \in A'$ such that
$\kappa(\mathfrak p) = \kappa(\mathfrak p')$,
$a_i - b_i \in (\mathfrak p')^N(A'_{\mathfrak p'})^\wedge$, and
$f_j(b_1, \ldots, b_n) = 0$ for $j = 1, \ldots, n$.
This determines an $R$-algebra map $B \to A'$ by sending the
class of $x_i$ to $b_i \in A'$. This finishes the proof
by taking $U = \Spec(A') \to \Spec(B)$ as the morphism $f$
and $u = \mathfrak p'$.
\end{proof}

\end{document}
